\section{Auditoría de fuentes y motivación: por qué se necesita una arquitectura de percepción general}

Esta clase no trata de «hacer otro Transformer de visión», sino de un problema de interfaz más fundamental: si la entrada puede provenir de imágenes, sonido, tacto, nubes de puntos, cámaras de eventos, texto, proteínas o instrumentos científicos, ¿es necesario rediseñar a mano, para cada tipo de datos, convoluciones, tokenizer, vecindades locales y decoder especializados? El objetivo de Perceiver es reducir al mínimo esta ingeniería modalidad por modalidad: primero trata los datos como un arreglo con características y posiciones, y luego completa la codificación, el procesamiento y la decodificación con una interfaz de attention unificada.

Esta reescritura toma como referencia temporal la carga oficial actual de Stanford Online `wTZ3o36lXoQ`. El video `GV8-6ZgJVRk` citado en las notas anteriores ya no está disponible; el video actual ofrece 1,399 subtítulos oficiales en inglés hechos a mano. La grabación no publica un PDF de slides independiente, por lo que se recuperaron 39 páginas didácticas completas a partir del video en 1080p; las animaciones repetidas, los saltos de página reiterados durante el Q\&A y los fotogramas intermedios de los videos de flujo óptico se fusionaron de forma deliberada, pero las explicaciones orales correspondientes se incorporaron al texto principal y al teacher-voice ledger.

\lecturefigure{slide-01-why-understand-all-modalities.jpg}{Por qué hay que entender todas las modalidades de datos: los datos del mundo real van mucho más allá de las imágenes y el texto.}{Video oficial actual de Stanford Online, 00:00:58--00:01:53.}

\begin{knowledgebox}{Qué es el inductive bias}
El inductive bias (sesgo inductivo) es una suposición estructural que la arquitectura incorpora de antemano. Por ejemplo, la convolución supone que la vecindad local es importante y que los pesos se comparten bajo traslación; el tokenizer supone ciertos límites discretos de subpalabras; las redes neuronales de grafos suponen que la entrada se compone de nodos y aristas. Los sesgos pueden mejorar la eficiencia en el uso de datos, pero también restringen las formas de datos a las que se aplican.
\end{knowledgebox}

\teachervoice{La primera razón de Drew Jaegle es muy pragmática: inventar a mano sesgos para cada tipo de sensor y de datos científicos simplemente no es algo que la comunidad de investigación pueda escalar de forma sostenida. El valor de una arquitectura general consiste, ante todo, en reducir el costo de «un esfuerzo de ingeniería de investigación por cada nueva modalidad».}

La segunda razón es el mantenimiento del sistema. Los enfoques multimodales tradicionales suelen construir por separado backbone, preprocesamiento y task head para video, audio y texto, y luego diseñar un módulo de fusión. Funcionan, pero a menudo solo sirven para unas pocas combinaciones fijas de entradas y son sensibles a la frecuencia de muestreo, a la forma de la cuadrícula y a la ausencia de alguna modalidad. Perceiver busca abstraer la construcción de la arquitectura para que los investigadores dediquen su esfuerzo a la tarea, los datos y el objetivo de aprendizaje.

\lecturefigure{slide-02-why-unified-systems.jpg}{De varios subsistemas especializados a un único modelo unificado capaz de recibir arreglos heterogéneos.}{Video oficial actual de Stanford Online, 00:02:58--00:05:09.}

\begin{importantbox}{«General» no significa «sin suposiciones»}
Perceiver sigue suponiendo que la entrada puede representarse como un arreglo finito, que attention puede aprender las relaciones entre elementos, que la información de posición o de modalidad puede proporcionarse como característica y que hay datos suficientes para recuperar la estructura de la tarea. Lo que reduce son las \textbf{suposiciones arquitectónicas específicas de cada modalidad}, no la eliminación de todo conocimiento previo.
\end{importantbox}

\begin{warningbox}{Una interfaz unificada y un entrenamiento unificado son dos cosas distintas}
La clase muestra la misma arquitectura aplicada por separado a varias tareas, pero la mayoría de estos experimentos se entrenaron de forma independiente. Poder reutilizar la interfaz encode/process/decode no equivale a haber entrenado un solo modelo que domine todas las modalidades a la vez.
\end{warningbox}

\subsection{Resumen de la sección}

El punto de partida de Perceiver es la expansión de modalidades y la complejidad de los sistemas. Busca una interfaz de arreglos reutilizable y menos estructura diseñada a mano; no afirma que el conocimiento del dominio, los modelos especializados o el entrenamiento multimodal conjunto hayan dejado de ser necesarios.
